% Source coverage: physical PDF pp. 31--35 (printed pp. 8--12), from the
% Appendix A heading through the final paragraph immediately before Appendix B.
% Inventory: Eqs. (24.A.1)--(24.A.7), one unnumbered display, no figures or
% tables, citations [2] and [14], and one source footnote.
% Corrected the printed cross-reference to the two supersymmetry-algebra sections.
\chapterappendix{A}{SU(6) Symmetry of Non-Relativistic Quark Models}

This appendix will describe the way that an $SU(6)$ symmetry that relates
particles of different spin arises in non-relativistic quark models. This has
nothing directly to do with supersymmetry, but it provides the historical
background for the Coleman--Mandula theorem, which is an essential input to the
construction of general supersymmetry algebras in Sections 25.2 and 32.1.

In general, the Hamiltonian of the non-relativistic quark model could depend not
only on positions and momenta but also on the spin and flavor operators
$\sigma_i^{(n)}$ and $\lambda_A^{(n)}$, where the $\sigma_i^{(n)}$ (with
$i=1,2,3$) act on the spin indices of the $n$th quark as the Pauli matrices
$\sigma_i$ defined by Eq.~(5.4.18), while the $\lambda_A^{(n)}$ (with
$A=1,2,\ldots,8$) act on the flavor indices of the $n$th quark as the
Gell-Mann $SU(3)$ matrices $\lambda_A$ defined by Eq.~(19.7.2). (Where $n$
refers to an antiquark, $\sigma_i^{(n)}$ and $\lambda_A^{(n)}$ act as the
matrices $-\sigma_i^{\mathsf T}$ and $-\lambda_A^{\mathsf T}$ of the contragredient
representations.) If we were to assume only that there is no spin-orbit
coupling, so that the total orbital angular momentum $L_i$ is separately
conserved, then we could conclude only that the Hamiltonian commutes with the
total spin and unitary spin
\begin{equation}
 S_i\equiv\frac{1}{2}\sum_n\sigma_i^{(n)},
 \qquad
 T_A\equiv\frac{1}{2}\sum_n\lambda_A^{(n)},
 \tag{24.A.1}\label{eq:24.A.1}
\end{equation}
as well as $L_i$. On the other hand, if we were to suppose that the Hamiltonian
depends only on quark positions and momenta, and not at all on spins or quark
flavors, then such a Hamiltonian would commute not only with the total orbital
angular momentum $L_i$, but also with \emph{each} of the operators
$\sigma_i^{(n)}$ and $\lambda_A^{(n)}$. In between these two extremes there is
the interesting possibility that in addition to commuting with the $L_i$,
$S_i$, and $T_A$, the Hamiltonian also commutes with the operators
\begin{equation}
 R_{iA}\equiv\frac{1}{2}\sum_n\mathord{\pm}
 \sigma_i^{(n)}\lambda_A^{(n)},
 \tag{24.A.2}\label{eq:24.A.2}
\end{equation}
where the sign is $+$ or $-$ for quarks and antiquarks,
respectively.\footnote{The minus sign for antiquarks arises because the terms
in $R_{iA}$ for antiquarks must act on spin and flavor indices as the matrix
$-(\sigma_i\lambda_A)^{\mathsf T}=-(-\sigma_i^{\mathsf T})(-\lambda_A^{\mathsf T})$.} The $S_i$, $T_A$, and
$R_{iA}$ form the Lie algebra of the group $SU(6)$, with commutation relations
\begin{equation}
\begin{aligned}
 [S_i,S_j]&=\ii\sum_k\epsilon_{ijk}S_k,
 &\qquad [T_A,T_B]&=\ii\sum_C f_{ABC}T_C,
 \\[2pt]
 [S_i,T_A]&=0,
 & [S_i,R_{jA}]&=\ii\sum_k\epsilon_{ijk}R_{kA},
 \\[2pt]
 [T_A,R_{iB}]&=\ii\sum_C f_{ABC}R_{iC},
 \\[2pt]
 [R_{Ai},R_{Bj}]&=\ii\delta_{ij}\sum_C f_{ABC}T_C
 +\frac{3}{2}\ii\delta_{AB}\sum_k\epsilon_{ijk}S_k
 \\[-1pt]
 &\quad+\ii\sum_{kC}\epsilon_{ijk}d_{ABC}R_{kC}.
\end{aligned}
\tag{24.A.3}\label{eq:24.A.3}
\end{equation}
Here $f_{ABC}$ and $d_{ABC}$ are respectively totally antisymmetric and totally
symmetric numerical coefficients~\hyperref[ch24-ref-14]{[14]}, with independent
non-vanishing values given by
\begin{equation}
\begin{aligned}
 f_{123}&=1,
 & f_{458}=f_{678}&=\sqrt{3}/2,
 \\
 f_{147}=f_{165}=f_{246}=f_{257}=f_{345}=f_{376}&=1/2,
\end{aligned}
\tag{24.A.4}\label{eq:24.A.4}
\end{equation}
and
\begin{equation}
\begin{aligned}
 d_{146}=d_{157}=-d_{247}=d_{256}=d_{344}=d_{355}
   =-d_{366}=-d_{377}&=1/2,
 \\
 d_{118}=d_{228}=d_{338}=-d_{888}&=1/\sqrt{3},
 \\
 d_{448}=d_{558}=d_{668}=d_{778}&=-1/(2\sqrt{3}).
\end{aligned}
\tag{24.A.5}\label{eq:24.A.5}
\end{equation}
This is the symmetry that remains if we include a spin- and flavor-dependent
two-body interaction in the Hamiltonian that commutes with $R_{iA}$ as well as
with the $S_i$ and $T_A$. There are such interactions, given by linear
combinations of two-body operators of the form
\begin{equation}
 H^{(nm)}\propto
 \left[1\mathbin{\pm}\sum_i\sigma_i^{(n)}\sigma_i^{(m)}\right]
 \left[\frac{2}{3}\mathbin{\pm}
 \sum_A\lambda_A^{(n)}\lambda_A^{(m)}\right],
 \tag{24.A.6}\label{eq:24.A.6}
\end{equation}
where the sign $\pm$ is negative if one of the particles $n$, $m$ is a quark
and the other an antiquark and positive if they are both quarks or both
antiquarks.

Of course, even in the non-relativistic quark model the $SU(6)$ symmetry is at
best approximate. It is broken by spin-orbit and spin-spin forces, and also by
the mass of the $s$ quark, which reduces the flavor $SU(3)$ symmetry to the
$SU(2)$ and $U(1)$ of isospin and hypercharge conservation. If we avoid the
effects of this quark mass difference by restricting ourselves to hadrons built
up from the light $u$ and $d$ quarks and antiquarks, then the only non-vanishing
$\lambda_A$ matrices are the $\lambda_a$ with $a=1,2,3$ (which for the $u$ and
$d$ quarks are given by the Pauli matrices (5.4.18) (that in this context are
conventionally called $\tau_a$) and $\lambda_8$ (which is just the number
$1/\sqrt{3}$ for the $u$ and $d$ quarks, and $-1/\sqrt{3}$ for the $\bar u$
and $\bar d$ antiquarks). The interaction \eqref{eq:24.A.6} thus becomes
\begin{equation}
 H^{(nm)}\propto
 \left[1\mathbin{\pm}\sum_i\sigma_i^{(n)}\sigma_i^{(m)}\right]
 \left[1\mathbin{\pm}\sum_A\tau_A^{(n)}\tau_A^{(m)}\right].
 \tag{24.A.7}\label{eq:24.A.7}
\end{equation}
Aside from quark number conservation, the remaining symmetry is then $SU(4)$,
with generators $S_i$, $T_a$, and $R_{ia}$ which commute with
\eqref{eq:24.A.7}. This was proposed by Wigner~\hyperref[ch24-ref-2]{[2]} in
1937 as a symmetry of nuclear forces, though of course with protons and
neutrons in place of $u$ and $d$ quarks. The interaction \eqref{eq:24.A.7} is
known in nuclear theory as a \emph{Majorana potential} to distinguish it from
an interaction that does not depend on spin or isospin, called a
\emph{Wigner potential}, or an interaction proportional to just the
spin-dependent or just the isospin-dependent factor in \eqref{eq:24.A.7},
known respectively as a \emph{Bartlett potential} and a
\emph{Heisenberg potential}.

It is amusing to note that, although in non-relativistic theories there is no
theoretical barrier to symmetries like $SU(6)$ that act on spin as well as
particle type, \emph{there never was any experimental evidence for such an
assumed $SU(6)$ symmetry of the non-relativistic quark model that is any better
satisfied than the assumption of complete spin and flavor independence.}
These assumptions are not the same; if the Hamiltonian of a system of $N$
non-relativistic quarks and/or antiquarks is completely independent of spin and
flavor, then its symmetry is $SU(6)^N$, not $SU(6)$. For instance, a
two-particle interaction like \eqref{eq:24.A.6} and various other multiparticle
interactions break $SU(6)^N$ to $SU(6)$. Of course, all these symmetries are
only approximate anyway. The question is whether $SU(6)$ is less badly broken
than $SU(6)^N$?

This cannot be answered by studying the multiplet that contains the baryon
octet, which consists of the nucleon and hyperons $\Lambda$, $\Sigma$, and
$\Xi$. In the non-relativistic quark model these particles are interpreted as
bound states of three quarks with zero orbital angular momentum. Because these
states are color neutral, the wave function is completely antisymmetric in the
suppressed color indices, and therefore it is completely symmetric under the
combined interchange of spin and flavor. The baryon octet would therefore have
to be put in the symmetric third-rank tensor representation $\mathbf{56}$ of
$SU(6)$, which besides the baryon octet contains a spin $3/2$ decuplet, which
may be identified as the one consisting of the famous `3--3' resonance
$\Delta$ and the $\Sigma(1385)$, $\Xi(1530)$, and $\Omega$ particles. (Numbers
in parentheses give masses in MeV, where these are needed to distinguish the
particles from others of the same isospin and strangeness but lower mass.) The
$SU(6)$ symmetry leads to good predictions for the baryon magnetic moments:
The quark charge operator is
$q=e(\lambda_3/2+\lambda_8/(2\sqrt{3}))$, so if quarks have the magnetic
moments $3q/(2m_N)$ of Dirac particles of this charge and mass $m_N/3$, then
the magnetic moment operator is
\[
 \mu_i\equiv 3\mu_N
 \left[\frac{1}{2}R_{i3}+\frac{1}{2\sqrt{3}}R_{i8}\right],
\]
where $\mu_N\equiv e/(2m_N)$ is the nuclear magneton, and $R_{iA}$ is defined
by Eq.~\eqref{eq:24.A.2}. It is straightforward to calculate the matrix
elements of this symmetry generator between the members of the $\mathbf{56}$
multiplet, with the result that the magnetic moments for $p$, $n$, $\Lambda$,
$\Sigma^+$, $\Sigma^-$, $\Xi^-$, and $\Xi^0$ in units of $\mu_N$ are
respectively $+3$, $-2$, $-1$, $+3$, $-1$, $-1$, and $-2$, which may be
compared with the corresponding experimental values $+2.79$, $-1.91$, $-0.61$,
$+2.46$, $-1.16$, $-0.65$, and $-1.25$. The agreement is fair, and somewhat
better (except for $\Sigma^-$) if we take the quark magnetic moment to be a
little less than $3\mu_N$. Because of the symmetry of the three-quark wave
function, nothing new is learned here if we assume that the Hamiltonian is
completely independent of spin and flavor; the states of zero angular momentum
would still have to fall in a multiplet consisting of
$6\cdot7\cdot8/6!=56$ members. In particular, the operator
\eqref{eq:24.A.6} has the same value 4 for any state of two quarks that is
symmetric under simultaneous interchange of their spins and flavors, whether
it is symmetric both under interchange of spins and interchange of flavors, or
antisymmetric under both interchanges.

In order to decide whether $SU(6)$ is any better than $SU(6)^N$, it is more
useful to study the mesons, which in the non-relativistic quark model are
interpreted as bound states of a quark and antiquark. If the Hamiltonian of
these states is completely independent of spin and flavor then its symmetry is
$SU(6)^2$, and the meson states fall into its 36-dimensional
$(\mathbf{6},\overline{\mathbf{6}})$ representation, while for $SU(6)$
symmetry we could only say that the mesons belong to either of the two
representations of $SU(6)$ contained in
$\mathbf{6}\times\overline{\mathbf{6}}$: the adjoint representation
$\mathbf{35}$ or the singlet representation. To be more specific, the
$\mathbf{35}$ consists of an $SU(3)$ singlet with spin $S=1$, an $SU(3)$
octet with $S=0$, and an $SU(3)$ octet with $S=1$, corresponding to the $SU(6)$
generators $S_i$, $T_A$, and $R_{iA}$, which is split from the $SU(3)$ singlet
state with $S=0$ by the interaction \eqref{eq:24.A.6}. Since all these
assumptions are approximate anyway, the question in deciding whether $SU(6)$
symmetry is more accurate than complete spin and flavor independence is
whether the splitting of the $SU(3)$ singlet $S=0$ state from the other
$\mathbf{35}$ states of the same orbital angular momentum is any greater than
the splittings within the $\mathbf{35}$ supermultiplet.

For orbital angular momentum $L=0$ the quark-antiquark states have negative
parity $P$, and positive or negative charge-conjugation quantum number $C$ (for
self-charge-conjugate states) according to whether the total spin $S$ is zero
or one, respectively. (For an explanation, see Section 5.5.) The $\mathbf{35}$
therefore consists of a singlet with $J^{PC}=1^{--}$, a $0^{-+}$ octet, and a
$1^{--}$ octet, which may be identified respectively as: the $\phi(1020)$; the
pseudoscalar octet $\pi$, $\eta$, $K$, and $\bar K$; and the vector octet
$\rho$, $\omega$, $K^*$, and $\bar K^*$. There is also a $0^{-+}$ $SU(3)$
singlet $\eta'$ at 958 MeV, which can be regarded as the $SU(6)$ singlet. The
splitting of this singlet from the particles in the $\mathbf{35}$ multiplet is
not distinctly greater than the splittings within the $\mathbf{35}$ multiplet.

It may be argued that the $L=0$ mesons do not provide a good test of the
symmetries of the non-relativistic quark model, because they include the
Goldstone bosons $\pi$, $\eta$, $K$, and $\bar K$, which become massless for
zero $u$ and $d$ quark masses and are therefore not well described by this
model. Therefore let us consider the quark-antiquark states with $L=1$. These
states have $P$ positive and $C$ positive or negative according to whether $S=1$
or $S=0$, so the p-wave $\mathbf{35}$ consists of: $SU(3)$ singlets with $S=1$
and $J^{PC}=0^{++}$, $1^{++}$, $2^{++}$, which may be identified as the
$f_0(1370)$, the $f_1(1285)$, and the $f_2(1270)$; an $S=0$ $1^{+-}$ octet
identified as $h_1(1170)$, $b_1(1235)$, $K_1(1400)$, and $\bar K_1(1400)$; and
$S=1$ octets: a $0^{++}$ octet consisting of $f_0(980)$, $a_0(980)$,
$K_0^+(1950)$, and $\bar K_0^+(1950)$; a $1^{++}$ octet consisting of
$f_1(1420)$, $a_1(1260)$, $K_1^+(1650)$, and $\bar K_1^+(1650)$; and a
$2^{++}$ octet consisting of $f_2(1430)$, $a_2(1320)$, $K_2^+(1980)$, and
$\bar K_2^+(1980)$. In addition to these $\mathbf{35}\cdot3$ states, there is
a another particle with the right quantum numbers to be the p-wave $SU(6)$
singlet: the $1^{+-}$ isoscalar $h_1(1380)$. Of course, we could interchange
the identifications of $h_1(1170)$ and $h_1(1380)$, or identify the $SU(3)$
singlet and octet isoscalar $1^{+-}$ states as orthogonal linear combinations
of $h_1(1170)$ and $h_1(1380)$. The important point is that there are
\emph{two} of these $1^{+-}$ isoscalars, with no indication that one of them,
belonging to an $SU(6)$ singlet, is more strongly split from the particles of
the $\mathbf{35}$ than the particles of the $\mathbf{35}$ are split from each
other. Here too, then, there is no evidence that $SU(6)$ symmetry is any more
accurate than the stronger assumption of complete spin and flavor independence.
